% =============================================================================
%  CHAPTER 20 -- PHASE 7: PORTFOLIO CONSTRUCTION AND RISK CONTROLS
% =============================================================================
\section{Phase 7 --- Portfolio-level construction and risk controls}
\label{ch:phase7}

\textbf{Reproduce it:} \code{python3 scripts/analysis/capture\_portfolio.py \&\& python3
scripts/plotting/render\_portfolio.py}.\\
\textbf{Math used:} Sections~\ref{sec:bt:risk} (VaR/CVaR, coherence) and
\ref{sec:bt:sizing} (inverse-vol, water-fill), \ref{sec:bt:regimes} (regime
conditioning).\\
\textbf{Artifacts:} \code{src/portfolio/risk.h},
\code{results/tables/portfolio\_\{summary,daily,regime,heat\_ez,heat\_ridge\}.csv},
\code{results/figures/portfolio\_\{equity,allocation,risk,heat\_ez,heat\_ridge\}.png},
\code{report/tables/portfolio.tex}.

\subsection{Goal}

A single basket is an anecdote. This phase builds a \emph{book}: four concurrent
market-neutral spread strategies, allocated by inverse volatility with enforced
concentration caps, measured with tail-risk statistics and decomposed by volatility
regime --- and then asks the question a fund would ask, which is not ``what is the
return?'' but ``what does the risk engineering buy?''.

\subsection{The four strategies}

Each is the identical Phase~\ref{ch:phase4} model --- causal 120-bar OLS hedge, $z$-bands
$2.0/0.5$, next-bar execution, 10\,bps one-way --- on a disjoint name set spanning three
GICS sectors:

\begin{center}
\renewcommand{\arraystretch}{1.2}
\begin{tabular}{llll}
\toprule
Target & Basket & Sector & Bars \\
\midrule
\code{ADBE} & \code{CRM}, \code{ADSK}, \code{INTU} & Information Technology (software) & 4{,}190 \\
\code{AVGO} & \code{ADI}, \code{INTC}, \code{AAPL} & Information Technology (semis/hardware) & 4{,}191 \\
\code{GS}   & \code{BAC} & Financials (large-cap banks) & 4{,}190 \\
\code{CAT}  & \code{HON} & Industrials (capital goods) & 4{,}190 \\
\bottomrule
\end{tabular}
\end{center}

The name sets are \emph{disjoint}, so the four daily P\&L streams are not mechanically
correlated and combining them is a real diversification statement rather than a
re-labelling of one position. Each strategy is dollar-neutral on its own
(equation~\eqref{eq:grossing}), so the aggregate book has near-zero net dollar exposure by
construction --- verified, not assumed, at every rebalance.

\subsection{Allocation: inverse volatility with a water-filled cap}

Weights are recomputed monthly (every 21 bars) from each strategy's trailing 252-day
realized daily P\&L volatility, $w_i\propto1/\sigma_i$ normalised to sum to one
(equation~\eqref{eq:inversevol}), then water-filled so that no strategy exceeds
\textbf{45\%} of book capital (Section~\ref{sec:bt:sizing}). Book-level limits: gross
weight $\le1.05$, absolute net weight $\le0.10$. Equal weighting is kept as the
comparison baseline.

At every rebalance the driver asserts and logs: $\sum_i w_i=1$ (book budget),
$w_i\le0.45$ (concentration), gross and net within limits. The primitives
(\code{inverse\_vol\_weights}, \code{gross\_exposure}, \code{net\_exposure},
\code{within\_limits}) are unit-tested in C++ independently of any market data
(\code{test\_inverse\_vol\_weights}, \code{test\_risk\_gross\_net\_limits}); the
water-fill itself lives in \code{capture\_portfolio.py::cap\_weights} and its
feasibility guard (if $\text{cap}\times m<1$ the constraint set is empty) reports
infeasibility rather than silently violating the cap.

\subsection{Results}

\input{tables/portfolio.tex}

\begin{figure}[H]\centering
\includegraphics[width=0.95\linewidth]{figures/portfolio_equity.png}
\caption{Combined NAV of the four-strategy book, net of costs, under equal-weight and
inverse-vol allocation, 2010--2026. Shaded bands mark the COVID-2020 and 2022-selloff
stress windows.}
\label{fig:portfolio_equity}
\end{figure}

\begin{figure}[H]\centering
\includegraphics[width=0.95\linewidth]{figures/portfolio_allocation.png}
\caption{Inverse-vol allocation across the four strategies, rebalanced monthly with a
45\% per-strategy cap (water-fill). The book budget $\sum w=1$ and the cap are verified
on every rebalance.}
\label{fig:portfolio_alloc}
\end{figure}

\begin{figure}[H]\centering
\includegraphics[width=0.95\linewidth]{figures/portfolio_risk.png}
\caption{Portfolio daily P\&L against the historical VaR$_{95}$ threshold (top) and the
log drawdown of the inverse-vol book (bottom). Exceedances cluster in the two stress
windows, as they should if the tail measure is measuring anything.}
\label{fig:portfolio_risk}
\end{figure}

Table~\ref{tab:portfolio_summary} is the generated record; the curated extract below is
the same data (net of 10\,bps, on \$1 gross):

\begin{center}
\renewcommand{\arraystretch}{1.2}
\begin{tabular}{llrrrrr}
\toprule
Book & Allocation & Net SR & Ann.\ net & MaxDD & VaR$_{95}$/day & CVaR$_{95}$/day \\
\midrule
\code{ADBE} & solo & $+0.050$ & $+0.38\%$ & $0.419$ & $0.00655$ & $0.01224$ \\
\code{AVGO} & solo & $-0.417$ & $-3.31\%$ & $0.777$ & $0.00704$ & $0.01297$ \\
\code{GS}   & solo & $-0.081$ & $-0.59\%$ & $0.312$ & $0.00693$ & $0.01200$ \\
\code{CAT}  & solo & $-0.437$ & $-4.52\%$ & $0.846$ & $0.00910$ & $0.01704$ \\
\midrule
Combined & equal weight & $-0.463$ & $-2.01\%$ & $0.440$ & $0.00409$ & $0.00681$ \\
Combined & \textbf{inverse vol} & $\mathbf{-0.401}$ & $\mathbf{-1.64\%}$ & $\mathbf{0.384}$ & $\mathbf{0.00389}$ & $\mathbf{0.00628}$ \\
\bottomrule
\end{tabular}
\end{center}

\subsection{The honest portfolio outcome}

\textbf{The book loses money.} Net Sharpe is $-0.401$ (inverse-vol) to $-0.463$ (equal
weight) over 2010--2026. Three of the four strategies are individually net-negative, and
the fourth (\code{ADBE}, $+0.050$) is not significant. This is the truthful result and it
is reported as the headline of the phase rather than buried under the risk statistics.

What the risk engineering demonstrably buys is \emph{tail reduction}, and it is measurable:
\begin{itemize}
  \item \textbf{Drawdown.} The worst solo books draw down $0.777$ (\code{AVGO}) and
    $0.846$ (\code{CAT}) of their gross notional; the combined book draws down $0.384$.
    Diversification cuts the worst path statistic by more than half, and inverse-vol cuts
    it a further $12.7\%$ versus equal weight ($0.440\to0.384$).
  \item \textbf{VaR and CVaR.} VaR$_{95}$ falls from $0.00409$ (equal) to $0.00389$
    (inverse-vol), $-4.9\%$; CVaR$_{95}$ falls from $0.00681$ to $0.00628$, $-7.8\%$.
    CVaR improving more than VaR is the signature of genuinely cutting the far tail rather
    than shifting the quantile --- and, per Section~\ref{sec:bt:risk}, CVaR is the coherent
    measure, so it is the one to trust.
  \item \textbf{Concentration discipline.} The 45\% cap prevents the allocation from
    piling into \code{ADBE} (the only net-positive book) simply because it had the quietest
    trailing year --- which is precisely the failure mode of uncapped inverse-vol
    weighting, since low recent volatility is not low future volatility.
\end{itemize}
What it does not buy is return: with every component's expected net return at or below
zero, no weighting of them produces a positive expectation. Risk controls manage the
distribution of a given edge; they cannot create one.

\subsection{Regime and stress decomposition}

\begin{itemize}
  \item \textbf{Volatility terciles} (1{,}376 days each, by trailing market volatility):
    net Sharpe $-0.183$ (low), $\mathbf{-0.781}$ (mid), $-0.344$ (high), with
    VaR$_{95}$ rising monotonically $0.00284\to0.00388\to0.00512$. The non-monotonicity
    in Sharpe is itself the interesting result: the middle tercile --- unsettled enough to
    break relationships, calm enough to keep trading them --- is the worst environment for
    this book, while the highest tercile at least stops trading often enough to limit the
    damage.
  \item \textbf{COVID-2020} (221 days): net Sharpe $-0.206$, VaR$_{95}$ $0.00555$.
    Bad, but not catastrophic --- the relationships bent rather than broke, and the
    dollar-neutral construction absorbed most of the market move.
  \item \textbf{2022 selloff} (251 days): net Sharpe $\mathbf{-1.306}$, VaR$_{95}$
    $0.00471$, annualised net return $-5.97\%$. By far the worst window: the
    rate-driven rotation out of growth equities broke exactly the software and
    semiconductor relationships that two of the four baskets are built on, and it broke
    them persistently enough that mean reversion kept paying the entry cost and never
    delivering the convergence.
  \item \textbf{2008 is not testable.} The panel begins 2010-01-04. The classic crisis for
    relative-value strategies (when correlations went to one and convergence trades
    diverged) is \emph{outside the sample} and cannot be measured. Disclosed here and in
    Chapter~\ref{ch:nongoals} rather than extrapolated.
\end{itemize}

\subsection{Parameter sensitivity}

\begin{figure}[H]\centering
\includegraphics[width=0.70\linewidth]{figures/portfolio_heat_ez.png}
\caption{Net Sharpe of the \code{ADBE} book as a function of regression lookback window
$W$ and entry threshold $z_{\text{entry}}$, at 10\,bps one-way.}
\label{fig:portfolio_heat_ez}
\end{figure}

\begin{figure}[H]\centering
\includegraphics[width=0.70\linewidth]{figures/portfolio_heat_ridge.png}
\caption{Net Sharpe of the \code{ADBE} book as a function of lookback window $W$ and ridge
penalty $\lambda$, at $z_{\text{entry}}=2.0$.}
\label{fig:portfolio_heat_ridge}
\end{figure}

The two surfaces (Tables \code{portfolio\_heat\_ez.csv}, \code{portfolio\_heat\_ridge.csv})
are the most actionable result of the whole project, and they say something consistent
with every cost calculation in Chapter~\ref{ch:backtestmath}:

\begin{center}
\renewcommand{\arraystretch}{1.2}
\begin{tabular}{lrrrr}
\toprule
$W$ & $z_{\text{entry}}=1.5$ & $2.0$ & $2.5$ & $3.0$ \\
\midrule
60  & $-0.244$ & $+0.076$ & $+0.044$ & $+0.153$ \\
120 & $-0.070$ & $+0.050$ & $+0.097$ & $+0.154$ \\
180 & $-0.165$ & $-0.081$ & $+0.137$ & $\mathbf{+0.206}$ \\
252 & $-0.117$ & $-0.033$ & $+0.087$ & $+0.151$ \\
\bottomrule
\end{tabular}
\end{center}

\begin{enumerate}[label=\textbf{P\arabic*.}]
  \item \textbf{Raising the entry threshold helps, monotonically and strongly.} At every
    window, net Sharpe rises with $z_{\text{entry}}$: at $W=180$,
    $-0.165\to-0.081\to+0.137\to+0.206$ as the threshold goes $1.5\to3.0$. The mechanism
    is \eqref{eq:costdrag}: a higher threshold trades less often, and each trade it does
    take starts from a larger expected reversion
    \eqref{eq:ou:profit}. Fewer, better-separated trades is the single biggest net-Sharpe
    lever available.
  \item \textbf{Longer windows help, but only with a high threshold.} The best cell
    overall is $(W{=}180, z{=}3.0)$ at $+0.206$ --- four times the baseline
    $(W{=}120, z{=}2.0)$ cell at $+0.050$. Longer windows give a smoother, less noisy
    hedge (the same mechanism as ridge in Phase~\ref{ch:phase5}); combined with a high
    entry bar, turnover collapses and the cost drag with it.
  \item \textbf{A short window with a low threshold is the worst configuration.}
    $(W{=}60, z{=}1.5)$ gives $-0.244$: a jumpy hedge and a trigger-happy policy is the
    maximum-turnover combination, and cost consumes everything.
  \item \textbf{Ridge helps at the windows where OLS churns.} On the $\lambda$ surface,
    $W{=}60$ improves from $+0.076$ ($\lambda{=}0$) to $+0.273$ ($\lambda{=}10$) and
    $+0.269$ ($\lambda{=}50$); $W{=}252$ improves from $-0.033$ to $+0.111$
    ($\lambda{=}10$) and $+0.173$ ($\lambda{=}50$). At $W{=}120$ the effect is small and
    non-monotone ($0.050\to0.041\to0.094\to0.025$). Shrinkage pays where the raw hedge is
    least stable --- exactly the variance-reduction story of
    equation~\eqref{eq:biasvariance}.
\end{enumerate}

\begin{remark}[Honesty: these surfaces are a search, and are priced as one]
The best cell ($+0.206$) was found by scanning $4\times4=16$ window/threshold
combinations and $4\times4=16$ window/$\lambda$ combinations on the \emph{same} data used
to report it. That is data snooping in its most ordinary form, and it is why every one of
these configurations is included in the $M=44$ audit of Phase~\ref{ch:phase10} rather than
being presented as a discovery. The defensible statement is the \emph{shape} of the
surface --- net Sharpe increases with entry threshold and with hedge smoothness, which is
a prediction of \eqref{eq:costdrag} and \eqref{eq:ou:profit} made before the scan --- not
the height of its best cell.
\end{remark}

\subsection{Exit gate --- status}

\begin{itemize}
  \item[\checkmark] multiple concurrent basket trades (four, disjoint names, three sectors);
  \item[\checkmark] inverse-vol sizing with per-strategy concentration cap, water-filled;
  \item[\checkmark] VaR/CVaR computed and compared across allocations;
  \item[\checkmark] gross/net limits enforced \emph{and verified} on every rebalance;
  \item[\checkmark] regime segmentation (vol terciles) and stress windows (2020, 2022) tabulated;
  \item[\checkmark] 2008 non-testability disclosed;
  \item[\checkmark] sensitivity heatmaps (lookback $\times$ threshold, lookback $\times$
    $\lambda$) generated from live runs.
\end{itemize}
